\section{Reproducibility, certification, and integration}
\label{sec:reproducibility}

\subsection{What the computations establish}

The proofs above establish the theorems; the computations independently
check normalization, signs, low-order specializations, and implementation.
They are deliberately recorded in three distinct categories.
\begin{enumerate}
\item \textbf{Exact algebra.} Partial-fraction identities are checked
as rational polynomial identities. Distribution and reflection matrices
are reduced over exact rational arithmetic. The modular differential
polynomials retain rational coefficients.
\item \textbf{Analytically certified enclosures.} The intervals in
Corollary~\ref{cor:rational} have exact rational endpoints. Their validity
comes from the proved spectral-tail bound, not a comparison of decimal
approximations.
\item \textbf{Independent high-precision checks.} Complex quadrature,
Hurwitz-zeta evaluation, and polygamma evaluation test the formulas by
different representations. These use ordinary arbitrary-precision
floating-point arithmetic, without outward rounding. They provide
strong implementation checks but are not interval certificates.
\end{enumerate}

The shipped default runs have the following scope.
\begin{center}
\small
\begin{tabular}{>{\raggedright\arraybackslash}p{0.27\textwidth}p{0.64\textwidth}}
\toprule
Artifact & Recorded verification\\
\midrule
Mixed doubles and Euler sums &
64 exact partial-fraction checks; 18 mixed integral comparisons;
the four explicit reductions; five Euler values; six generator checks;
three pole-residue checks. Working precision: 80 decimal digits.\\
\addlinespace
Distributions and bridges &
1,734 exact matrix checks through \(q=60\), including anchored and
reflected quotients and all-divisor/prime-row agreement.
72 analytic checks of normalized bridges through derivative order five,
at 65 decimal digits.\\
\addlinespace
Parameter asymptotics &
78 rational enclosures, including 72 strictly positive sign certificates;
72 independent exact coefficient comparisons; 36 numerical
spectral-tail checks at 75 decimal digits; the profile figure.\\
\addlinespace
Modular rows &
Four independent single-row comparisons, ten CM jet comparisons through
order four, and four explicit tail checks at 90 decimal digits;
exact differential polynomials through order five.\\
\bottomrule
\end{tabular}
\end{center}

The maximum equality residual in the mixed/Euler run was
approximately \(1.71\times10^{-80}\), against the declared tolerance
\(10^{-60}\). Pole residues use their own finite-difference error
criterion and are not included in that equality statistic.
The source manifest verifies all 40 retrieved text files against their
Git blob identifiers and records additional SHA-256 hashes.
This is provenance for the reviewed snapshot, not a claim that every
statement in all 31 reports has been checked.

\subsection{A certificate format that survives integration}

The exact objects should remain separate from their numerical
evaluations. For a linear relation, store the convention and ordered
symbol list, the integer or rational coefficient vector, the source
law identifiers, and the exact combination of law rows that produces
the vector. An integer-relation search may propose that vector, but
does not supply its law certificate.

For a finite-dimensional quotient, store the relation system and an
exact normal-form map in addition to the rank. A rank number without
the rows and their interpretation is not enough to reconstruct the
claim. The character basis in Section~\ref{sec:distribution} provides
a particularly compact symbolic normal form across infinitely many
denominators. Its general proof is stronger than a larger finite table.

For numerical values, store the representation, parameters, precision,
truncation, and the actual error statement. Where a proved analytic
majorant is evaluated in ordinary floating point, label the result
accordingly; an outward-rounded evaluation is an additional step.
For the rational intervals, retain numerator and denominator rather
than only a decimal rendering.

The package is self-contained at the source level: the master
\src{article.tex} includes the section files, the bibliography is
included, the figure source and data are present, and a build command
is supplied in \src{README.txt}. The Python scripts use
\src{mpmath}, \src{sympy}, \src{numpy}, and \src{matplotlib}.
No absent upstream value store is needed to reproduce these results.
The proposed corrections are supplied as a separate register so that
maintainers can review and apply them to the original articles.
The original corpus is neither duplicated nor silently modified.

\section{Further research: precise questions and proposed routes}
\label{sec:questions}

The corrections change which problems are worth pursuing. In
particular, the inverse-argument mixed family no longer supplies the
proposed new directions. A productive continuation should use exact
relations to remove those directions before allocating numerical
effort to the genuinely unresolved quotient.

\subsection{Immediate structural projects}

\paragraph{1. A convention-safe cyclotomic reduction engine.}
Can the partial-fraction map of Theorem~\ref{thm:mixed} be combined
with shuffle, stuffle, distribution, conjugation, and the proven
parity theorem into a single exact reducer at fixed weight and level?
Start with levels three, four, and six, retain an explicit ordered
alphabet, and distinguish products from the linear depth filtration.
Every emitted relation should carry a rational certificate.
The first target is a reproducible spanning set and formal corank
through a stated weight, not an arithmetic dimension theorem.
The four false survivors make useful regression examples.

\paragraph{2. The remaining parity of alternating harmonic sums.}
For
\[
S_{2m}=\sum_{n\ge0}\frac{(-1)^nH_n}{(2n+1)^{2m}},
\]
what is the smallest natural motivic or formal level-four space that
contains the family? The beta-integral argument evaluates the other
parity but does not prove nonmembership here. A concrete next step is
to compute a depth-graded coaction or an exact formal quotient for the
first few weights, identify the image of \(S_{2m}\), and distinguish
the resulting motivic statement from any period-independence
assumption. Xu and Wang's wider Euler-sum framework is a natural
starting point~\cite{XuWang}.

\paragraph{3. Higher-depth inverse-color reductions.}
The substitution \(n=m+h\) works especially well because the inverse
arguments leave a single phase \(z^h\). For triples such as
\(\Li_{a,b,c}(z,z^{-1},1)\), which phase patterns lead to a similar
gap-variable reduction, and which genuinely retain two interacting
gaps? A first target is an exact partial-fraction identity at depth
three with a complete boundary-convergence proof. Merely increasing
the depth-two search basis cannot establish that a triple lies in it.

\paragraph{4. Integral weighted distributions.}
The character proof works over a \(\C\)-algebra, where averaging over
\((\Z/q\Z)^\times\) and diagonalizing characters are available.
What is the corresponding integral module over
\(\Z[t_p:p\mid q]\), and which torsion appears after reflection,
anchoring, or specialization? The distinction matters if a future
certificate system is to use integral normal forms without inverting
group orders. Smith normal forms for prime powers and two-prime
denominators would supply useful conjectures; an induction or a
distribution resolution is needed for the general theorem.
The connection to universal distributions should be retained
\cite{Kubert,Ouyang}.

\subsection{Asymptotic and analytic projects}

\paragraph{5. Beyond the logarithmic index range.}
Theorem~\ref{thm:profile} is uniform for \(n\le C\log(k/a)\).
What happens when \(n/\log k\to\infty\)? The exact coefficient
representation suggests a saddle-point equation for
\[
e^{-z\log a}\prod_{j=1}^{k}(1+z/j),
\]
but the appropriate saddle may move far enough that fixed-disk
gamma-ratio estimates no longer apply. The competing \(m\ge1\)
spectral terms also need uniform control.
A concrete first extension is a proved range
\(n=o((\log k)^2)\), with an error estimate that explicitly describes
where the two-term profile ceases to be useful. That range is a
research target here, not an established theorem.

\paragraph{6. Moving Hurwitz parameter and spectral competition.}
The present compact-\(a\) hypothesis gives
\((a/(a+1))^k\) as an exponentially small ratio.
For \(a=a(k)\) comparable with \(k\), this ratio need not tend to zero.
Can one derive a profile involving several spectral terms, or an
integral limit of them? The separate regimes
\(a=o(k)\), \(a\asymp k\), and \(a\gg k\) should be analyzed using
the exact expansion before proposing a single universal formula.
This would connect the parameter tower with large-parameter
Hurwitz-zeta asymptotics.

\paragraph{7. Fully effective asymptotic constants and optimal bounds.}
Lemma~\ref{lem:coefficient} has explicit constants, while
Theorem~\ref{thm:profile} packages the fixed-disk bounds into \(O\)-terms.
An effective version could expose a computable \(k_0(a_0,a_1,C)\)
for the sign theorem. A separate optimization problem is to minimize
\eqref{eq:spectral-tail} over \(R\) and the truncation \(M\), or to
compare the integer-radius rational certificate with a rational
upper approximation to the optimal real radius.
The success criterion is a rigorously narrower interval at a known
computational cost, not simply more printed digits.

\paragraph{8. Resonant jet coordinates and numerical conditioning.}
The full distribution module is free even where top-denominator
coordinates vanish. How should a symbolic or numerical algorithm
choose conductor-level coordinates so that the conversion remains
well-conditioned near a resonance? The vanishing-order formula in
Section~\ref{sec:distribution} provides a local answer, including the
number of lost jet coefficients. A useful next deliverable is a
canonical basis-selection algorithm with explicit transition matrices,
followed by a study of conditioning at large denominators.
No extra numerical constants should be inferred from a coordinate
singularity.

\paragraph{9. Global continuation of normalized derivative bridges.}
Section~\ref{sec:bridges} gives a local germ after removing a simple
trivial zero. How far can each germ be continued before an actual
zero of the \(L\)-function or a known elementary factor becomes an
obstruction? One can formulate zero-sensitive domains for the
logarithmic jets, derive Bell-polynomial formulas for raw derivatives,
and build exact symbolic identities for higher orders.
For imprimitive characters, Euler-factor zeros should be removed
explicitly and compared with the distribution resonances.
The root-number cancellation in the local identity should not be
misread as a global zero-free theorem.

\subsection{Arithmetic and modular questions}

\paragraph{10. The unresolved arithmetic content of gamma quotients.}
The formal gamma quotient has a proved rank, and every law-derived
identity has an exact certificate. Which additional arithmetic
relations could remain beyond that system? This is the point at
which Rohrlich-type conjectures enter. A useful project would separate
three outputs: unconditional law certificates, conditional
completeness statements with the exact conjecture written out, and
finite bounded-height exclusions using certified numerical inputs.
The Koblitz--Ogus sufficient criterion must not be promoted to the
missing converse; see the careful distinction in
Anderson--Brownawell--Papanikolas~\cite{ABP}.

\paragraph{11. Isolated imaginary polygamma values at CM points.}
The modular row theorems evaluate weighted aggregates. They do not
individually evaluate \(\Impart\psi^{(1)}(\ii)\).
Does this isolated quantity belong to a natural Eichler-integral,
iterated-integral, or regulator framework with a provable modular
transformation law? The immediate target is the transformation law
itself, including its nonholomorphic or boundary terms if necessary.
A Lambert-series resemblance alone is not sufficient evidence.
Once such a law is established, it may suggest period classes and
meaningful relations among different CM arguments.

\paragraph{12. Beyond the square and hexagonal lattices.}
At other imaginary-quadratic parameters, the same Ramanujan
recurrence computes all row derivatives from \(E_2,E_4,E_6\).
Can the algebraic CM data and periods be certified exactly, with
minimal polynomials, isolating regions, and compatible period
normalizations? A useful first target is several class-number-two
discriminants. A polynomial relation supplies an upper bound on
degree; exact degree additionally requires irreducibility and
identification of the intended root. This would replace the draft's
PSLQ degree guesses by checkable algebraic evidence.

\subsection{A sensible order of attack}

The first, fourth, seventh, and eighth projects have the clearest
near-term certification goals. They strengthen the mathematical
infrastructure without depending on a period conjecture. The second,
third, and ninth concern deeper structural reductions, where formal
or motivic calculations can establish substantial results before
arithmetic independence is addressed. The remaining asymptotic and CM
projects require new analytic input.

The central methodological requirement is to keep three maps visible:
from exact functional identities to a formal quotient, from that
quotient to actual special values, and from those values to numerical
approximations. A proof about one map should not silently become a
claim about another.
